\chapter{Opportunity Identification in Hydroponic Agriculture}

\section{Session Overview}

This session begins the project by asking your group to identify a realistic product opportunity within hydroponic agriculture and to place that opportunity inside a believable company context. These decisions matter because later activities, such as customer-needs analysis, specification setting, and concept generation, depend on having a clear product scope from the start.

By the end of this session, your group should be able to:
\begin{enumerate}
    \item select a promising product opportunity within hydroponic agriculture;
    \item justify that choice using market need, feasibility, user impact, and team capability;
    \item choose a suitable company context for the project; and
    \item document the decision in a concise point-form summary for later sessions.
\end{enumerate}

\section{Hydroponics as the Project Context}

Hydroponics is a method of growing plants without soil by controlling water, nutrients, support structures, and environmental conditions. For Industrial Physics students, this theme is useful because many hydroponic products depend on measurable technical behaviour such as fluid flow, structural loading, sensor response, power use, heat transfer, and control-system performance.

A strong Session 1 idea should therefore satisfy two conditions:
\begin{itemize}
    \item it should solve a meaningful hydroponics-related problem; and
    \item it should create enough technical depth for later engineering analysis.
\end{itemize}

\section{What Counts as a Suitable Product?}

This section defines the types of products your group may consider.

A \textbf{direct product} performs a core hydroponic function inside the growing system. Examples include a nutrient-mixing device, a water-distribution component, a plant-support module, or a sensing unit that monitors pH or electrical conductivity.

An \textbf{indirect product} supports hydroponic work without becoming the main growing mechanism itself. Examples include a maintenance tool, a modular storage unit for hydroponic supplies, or a harvesting aid designed for hydroponic users.

In both cases, the link to hydroponics must be specific and defensible. Avoid ideas that are too broad, too generic, or only loosely related to the project theme.

\section{Example Opportunity Areas}

Table~\ref{tab:session1-opportunity-areas} gives example product areas. These examples are intended to guide discussion, not to limit your group to a fixed list. Use the table to identify which opportunities are most suitable for your group's interests and technical strengths.
\begin{table}[htbp]
\centering
\caption{Example product areas in hydroponic agriculture and related physics concepts}
\label{tab:session1-opportunity-areas}
\renewcommand{\arraystretch}{1.2}
\begin{tabularx}{\textwidth}{|>{\raggedright\arraybackslash}p{0.22\textwidth}|>{\raggedright\arraybackslash}X|>{\raggedright\arraybackslash}p{0.28\textwidth}|}
\hline
\textbf{Product area} & \textbf{Example opportunities} & \textbf{Relevant physics concepts} \\
\hline
Seedling and germination & Seed trays, seed holders, nursery supports, moisture-control inserts & Heat and mass transfer, moisture control, material properties \\
\hline
Nutrient preparation and delivery & Mixing containers, dosing tools, pumps, tubing connectors, flow regulators & Fluid dynamics, pressure loss, pump performance, mixing behaviour \\
\hline
Plant support and structure & Net pots, support frames, vertical racks, modular channels & Structural mechanics, load distribution, material stiffness \\
\hline
Monitoring and control & pH indicators, water-level monitors, sensor housings, lighting supports & Instrumentation, signal response, calibration, control systems \\
\hline
Maintenance and cleaning & Drainage tools, cleaning devices, inspection aids, quick-replacement parts & Flow management, ergonomics, corrosion resistance, durability \\
\hline
Harvesting and handling & Harvest trays, cutting aids, storage bins, packaging supports & Force, safety, material selection, impact protection \\
\hline
\end{tabularx}
\end{table}

\section{Activity 1: Select the Product Opportunity}

This activity asks your group to choose one product to carry into later sessions. Start by brainstorming several candidates, then compare them using the criteria below.

\subsection{Evaluation Criteria}

When comparing product ideas, ask the following questions:
\begin{enumerate}
    \item \textbf{Need}: Does the product address a real hydroponics problem or user frustration?
    \item \textbf{Feasibility}: Can your group study and develop the idea within the time and resource limits of this course?
    \item \textbf{User impact}: Would the product make operation, maintenance, monitoring, or harvesting meaningfully better?
    \item \textbf{Physics relevance}: Does the idea allow you to apply industrial physics concepts in a clear way?
    \item \textbf{Team capability}: Does your group have enough interest and confidence to work on the idea for the full semester?
\end{enumerate}

Your product may be a new idea or an improved version of an existing market product. At this stage, you only need a quick market scan to understand whether similar products already exist and how your group might position the concept.

\subsection{Expected Output}

At the end of this stage, your group should be able to state:
\begin{itemize}
    \item the selected product;
    \item the specific hydroponics problem it addresses; and
    \item the reasons it was preferred over other candidate ideas.
\end{itemize}

\section{Activity 2: Select the Company Context}

After choosing the product, decide which company your group will represent. This step matters because the company context influences later choices about market segment, manufacturing capability, pricing, features, and design ambition.

Your group may choose either of the following company types:
\begin{itemize}
    \item an \textbf{established company} that already sells related agricultural or technical products; or
    \item a \textbf{start-up or small company} entering the market with a focused offer.
\end{itemize}

For example, a modular hydroponic growing rack could be developed either by an established agricultural-equipment manufacturer expanding its product line or by a start-up targeting urban growers who need compact systems.

When choosing the company context, discuss:
\begin{enumerate}
    \item the likely target users;
    \item the resources the company could realistically access;
    \item the quality level and market position of the product; and
    \item whether the company context supports your proposed selling price and production approach.
\end{enumerate}

\section{Optional Extension Work}

The tasks in this section are not required to be completed during class, but they are useful if your group wants a stronger starting point for future sessions. Divide these tasks according to team interest and expertise.

\begin{enumerate}
    \item \textbf{Explore customer needs}: identify likely user frustrations, desired improvements, and contexts of use.
    \item \textbf{Sketch the design direction}: outline the main functions, components, and performance considerations.
    \item \textbf{Estimate product structure}: list possible parts, interfaces, or modules and sketch a rough arrangement.
    \item \textbf{Estimate market position}: propose an approximate selling price and note the assumptions behind it.
\end{enumerate}

If your group chooses to go further, assign responsibilities explicitly. For example, one member may scan the market, another may review technical feasibility, and another may sketch early product ideas.

\section{Next Steps}

Before leaving this session, compile a short point-form summary containing:
\begin{itemize}
    \item the chosen product;
    \item its relationship to hydroponic agriculture;
    \item the reasons for selecting it;
    \item the chosen company context; and
    \item the rationale for that company choice.
\end{itemize}

Bring this summary to the next session and upload it to ITEL. A clear decision here will make later work on planning, customer needs, and specifications much easier.
